\documentclass[11pt]{article}
\usepackage[margin=1in]{geometry}
\usepackage{graphicx}
\usepackage{booktabs}
\usepackage{amsmath}
\usepackage{hyperref}
\usepackage{caption}
\usepackage{subcaption}
\usepackage{xcolor}
% multirow not available; use manual spacing instead

\hypersetup{colorlinks=true, linkcolor=blue, citecolor=blue, urlcolor=blue}

\title{Widest-Path vs.\ Shortest-Path Routing\\Under CSMA/Bianchi WiFi Interference}
\author{ncsim Experiment Report}
\date{\today}

\begin{document}
\maketitle

\begin{abstract}
We compare \textbf{widest-path} and \textbf{shortest-path} multi-hop routing
strategies for DAG-scheduled task graphs on wireless mesh networks using the
ncsim discrete-event simulator with its expected-rate CSMA/Bianchi interference
model and HEFT scheduling. On \textbf{grid mesh topologies} (9 scenarios:
$2\times2$, $3\times3$, $4\times4$ grids $\times$ 5, 10, 20 tasks),
shortest-path wins 7 of 9 cases (2 ties), with widest-path slowdowns up to
$4.4\times$ on larger networks. The root cause is \emph{flow concentration}:
widest-path funnels transfers onto high-bandwidth links, creating CSMA
contention. However, on a \textbf{bottleneck relay network} where a slow direct
link ($1.075$\,MB/s) competes with a high-bandwidth relay chain
($8.6$\,MB/s per hop), widest-path wins all 3 scenarios with speedups up to
$2.5\times$. The results demonstrate that optimal routing depends on network
structure: shortest-path is preferred for dense meshes, while widest-path
excels when relay paths offer substantially higher bandwidth than direct links.
\end{abstract}

%% ============================================================
\section{Introduction}

When scheduling DAG-structured computations across a wireless mesh network,
data must be transferred between nodes after each task completes. The
\emph{routing strategy} determines which multi-hop path each transfer takes.
Two natural choices are:

\begin{itemize}
    \item \textbf{Widest path}: maximize the minimum link bandwidth along the
    route (Dijkstra-style, maximizing bottleneck capacity).
    \item \textbf{Shortest path}: minimize the number of hops (BFS-style,
    minimizing hop count).
\end{itemize}

In isolation, widest-path routing selects higher-throughput links. However, in
a wireless network where all links within carrier-sensing range contend for
airtime via CSMA, concentrating flows on fewer ``best'' links can create
severe contention. Conversely, when the shortest path traverses a slow
bottleneck link while an alternative relay chain offers much higher bandwidth,
widest-path routing can dramatically outperform shortest-path. This report
quantifies both effects using ncsim's expected-rate CSMA/Bianchi interference
model across grid mesh and bottleneck relay topologies.

%% ============================================================
\section{Experimental Setup}

\subsection{Network Topologies}

We use square grid mesh networks with 40\,m spacing between adjacent nodes.
Links include all grid edges (horizontal, vertical) and checkerboard-pattern
diagonals, all bidirectional. The WiFi PHY rates are computed from the
802.11ax MCS table based on SNR at each link distance: 40\,m grid links and
$\approx$56.6\,m diagonal links receive different PHY rates, giving widest-path
and shortest-path different preferred routes.

\begin{table}[h]
\centering
\caption{Network configurations.}
\label{tab:networks}
\begin{tabular}{lccc}
\toprule
Label & Grid & Nodes & Bidirectional Links \\
\midrule
Small   & $2\times2$ & 4  & 12 \\
Medium  & $3\times3$ & 9  & 32 \\
Large   & $4\times4$ & 16 & 68 \\
\bottomrule
\end{tabular}
\end{table}

Nodes have heterogeneous compute capacities (80--300 compute\_units/s) so that
HEFT distributes tasks across multiple nodes, creating inter-node data
transfers.

\subsection{DAG Structures}

Three DAG sizes capture different communication patterns:

\begin{table}[h]
\centering
\caption{DAG configurations (all tasks: compute\_cost=500, all edges: data\_size=10\,MB).}
\label{tab:dags}
\begin{tabular}{lcl}
\toprule
Label & Tasks & Shape \\
\midrule
Small  & 5  & Fork-join: 1 source $\to$ 3 parallel $\to$ 1 sink \\
Medium & 10 & Diamond: source $\to$ 4 $\to$ 4 $\to$ sink \\
Large  & 20 & Multi-level: 5-stage pipeline with branching \\
\bottomrule
\end{tabular}
\end{table}

\subsection{Simulation Parameters}

All simulations use:
\begin{itemize}
    \item \textbf{Scheduler}: HEFT (Heterogeneous Earliest Finish Time)
    \item \textbf{Interference}: CSMA/Bianchi (802.11ax, 5\,GHz, 20\,dBm TX, path loss exponent 3.0)
    \item \textbf{Carrier sensing range}: 71.2\,m (all links within each grid are in the conflict graph)
    \item \textbf{Seed}: 42 (deterministic)
\end{itemize}

The CSMA/Bianchi model combines two interference mechanisms: (1)~\emph{Bianchi
contention} for links sharing a conflict-graph edge (symmetric MAC-layer
time-sharing), and (2)~\emph{SINR degradation} for hidden terminals outside
the conflict graph. In our grid topologies, nearly all links are within
carrier-sensing range, so Bianchi contention dominates.

%% ============================================================
\section{Aggregate Results}

\begin{table}[h]
\centering
\caption{Makespan (seconds) for all 18 grid simulations. Bold indicates the winner in each cell. Percentage shows widest-path overhead relative to shortest-path.}
\label{tab:results}
\begin{tabular}{ll rrr}
\toprule
& & \multicolumn{3}{c}{Network Size} \\
\cmidrule(lr){3-5}
DAG & Routing & $2\times2$ (4) & $3\times3$ (9) & $4\times4$ (16) \\
\midrule
Small (5)
  & Widest   & \textbf{11.43} & \textbf{12.37} & 39.15 \\
  & Shortest & \textbf{11.43} & \textbf{12.37} & \textbf{8.93} \\
  & \emph{Diff} & \emph{tie} & \emph{tie} & \emph{+338\%} \\
\midrule
Medium (10)
  & Widest   & 30.02          & 42.76          & 124.86 \\
  & Shortest & \textbf{20.52} & \textbf{32.33} & \textbf{101.62} \\
  & \emph{Diff} & \emph{+46\%} & \emph{+32\%} & \emph{+23\%} \\
\midrule
Large (20)
  & Widest   &  51.65         &  70.19         & 221.73 \\
  & Shortest & \textbf{36.66} & \textbf{51.36} & \textbf{160.62}  \\
  & \emph{Diff} & \emph{+41\%} & \emph{+37\%} & \emph{+38\%} \\
\bottomrule
\end{tabular}
\end{table}

\begin{table}[h]
\centering
\caption{Winner summary for grid topologies: shortest-path wins in 7 of 9 cases.}
\label{tab:winners}
\begin{tabular}{lccc}
\toprule
& $2\times2$ & $3\times3$ & $4\times4$ \\
\midrule
Small DAG  & Tie      & Tie      & Shortest \\
Medium DAG & Shortest & Shortest & Shortest \\
Large DAG  & Shortest & Shortest & Shortest \\
\bottomrule
\end{tabular}
\end{table}

The dominant trend is clear: shortest-path routing wins in all non-trivial
grid scenarios. The two ties occur with the small DAG on small networks, where
both strategies use the same single-hop path. The widest-path overhead ranges
from 23\% ($4\times4$ medium DAG) to 338\% ($4\times4$ small DAG), with the
largest absolute slowdown on the $4\times4$ grid where widest-path routes
transfers through long perimeter corridors.

%% ============================================================
\section{Per-Scenario Analysis}

Each figure below shows the widest-path (left) and shortest-path (right)
routing side by side. The top row displays the network topology with links
colored and thickened by flow count (number of data transfers using that link).
The middle row shows a Gantt-style timeline of task executions (blue) and
data transfers (orange). The bottom row shows per-directed-link flow counts.

%% --- 2x2 Small ---
\clearpage
\subsection{$2\times2$ Network $\times$ Small DAG (5 tasks)}

\begin{figure}[h]
\centering
\includegraphics[width=\textwidth]{figures/smallnet_smalldag.png}
\caption{$2\times2$ $\times$ Small DAG. Both routing strategies produce identical results.}
\label{fig:ss}
\end{figure}

With only 4 nodes and a simple fork-join DAG, HEFT places 4 of 5 tasks on the
fastest node (n0, capacity=200) and sends one task to n2. Both strategies use
the same single-hop direct link (n0$\leftrightarrow$n2) for both transfers.
There is no routing decision to make: the network is too small for path
diversity.

%% --- 2x2 Medium ---
\subsection{$2\times2$ Network $\times$ Medium DAG (10 tasks)}

\begin{figure}[h]
\centering
\includegraphics[width=\textwidth]{figures/smallnet_mediumdag.png}
\caption{$2\times2$ $\times$ Medium DAG. Shortest path wins (20.52\,s vs.\ 30.02\,s).}
\label{fig:sm}
\end{figure}

With 10 tasks across 3 nodes, widest path introduces 2-hop routes through
diagonal links, increasing total link occupancy and CSMA contention. Under the
expected-rate SINR model, the additional hidden-terminal interference from
multi-hop routes outweighs any bandwidth advantage. Shortest path keeps all
transfers on direct 1-hop links, minimizing the contention footprint and
finishing 46\% faster.

%% --- 2x2 Large ---
\clearpage
\subsection{$2\times2$ Network $\times$ Large DAG (20 tasks)}

\begin{figure}[h]
\centering
\includegraphics[width=\textwidth]{figures/smallnet_largedag.png}
\caption{$2\times2$ $\times$ Large DAG. Widest path concentrates 14 flows on link n0--n2.}
\label{fig:sl}
\end{figure}

With 20 tasks on 3 nodes, the flow-concentration problem is pronounced. Widest
path introduces 2-hop routes, doubling link occupancy per transfer, and
concentrates flows onto high-bandwidth diagonal links. Shortest path keeps all
transfers on direct 1-hop links, distributing traffic more evenly across the
available physical links. The Gantt chart shows widest path's transfers
stretching to $\sim$49\,s while shortest path finishes by 36.7\,s (41\%
improvement).

%% --- 3x3 Small ---
\subsection{$3\times3$ Network $\times$ Small DAG (5 tasks)}

\begin{figure}[h]
\centering
\includegraphics[width=\textwidth]{figures/mediumnet_smalldag.png}
\caption{$3\times3$ $\times$ Small DAG. Both strategies take the same 2-hop path through n3.}
\label{fig:ms}
\end{figure}

HEFT concentrates 4 of 5 tasks on n4 (center, capacity=300) and places one
task on a corner node. Both strategies must traverse the same 2-hop path
(n4$\to$n3$\to$n6 and back). With only 2 sequential transfers and no
alternative paths, routing is irrelevant.

%% --- 3x3 Medium ---
\clearpage
\subsection{$3\times3$ Network $\times$ Medium DAG (10 tasks)}

\begin{figure}[h]
\centering
\includegraphics[width=\textwidth]{figures/mediumnet_mediumdag.png}
\caption{$3\times3$ $\times$ Medium DAG. Widest path scatters traffic through n0 and n6
via 10 different links; shortest path uses 6 links concentrated around n4's neighbors.}
\label{fig:mm}
\end{figure}

The topology panels diverge clearly. Widest path routes traffic through
distant nodes (n0 and n6) via 2-hop paths through n1 and n3, activating
\textbf{10 unique links}. Shortest path concentrates on n4's direct neighbors
(n7 and n6) using only \textbf{6 unique links} with a lower average hop count
(1.5 vs.\ 2.0). Widest path's longer routes mean more simultaneous link
activations, increasing CSMA/Bianchi contention across the conflict graph.
Shortest path finishes 32\% faster (32.33\,s vs.\ 42.76\,s).

%% --- 3x3 Large ---
\subsection{$3\times3$ Network $\times$ Large DAG (20 tasks)}

\begin{figure}[h]
\centering
\includegraphics[width=\textwidth]{figures/mediumnet_largedag.png}
\caption{$3\times3$ $\times$ Large DAG. Widest path creates a congested n4--n3--n6 corridor
with 6 flows per link; some transfers include 3-hop routes.}
\label{fig:ml}
\end{figure}

With 20 tasks distributed across 4 nodes, widest path creates a heavily
loaded corridor: multiple flows on n4$\to$n3 and n3$\to$n6. The topology panel
shows thick red links forming an L-shaped corridor on the left side of the
grid. Widest path also uses longer routes (avg 2.0 hops, max 3), increasing
per-transfer link occupancy. Shortest path uses a similar number of unique
links but with lower hop counts (avg 1.6, max 2), completing 37\% faster
(51.36\,s vs.\ 70.19\,s).

%% --- 4x4 Small ---
\clearpage
\subsection{$4\times4$ Network $\times$ Small DAG (5 tasks)}

\begin{figure}[h]
\centering
\includegraphics[width=\textwidth]{figures/largenet_smalldag.png}
\caption{$4\times4$ $\times$ Small DAG. Widest path takes an 8-hop perimeter route to n12;
shortest path uses a 1-hop direct link to n0. The $35\times$ slowdown is the most extreme case.}
\label{fig:ls}
\end{figure}

This is the most dramatic grid case. HEFT assigns one task to a remote node
via widest-path routing, traversing a multi-hop perimeter route. Each hop
adds latency, and with CSMA/Bianchi, all links in the route must contend for
airtime simultaneously. In contrast, shortest path sends the task to the
adjacent node n0 via a \textbf{single 1-hop link}, completing much faster. The
makespan ratio is $39.15 / 8.93 = 4.4\times$.

%% --- 4x4 Medium ---
\subsection{$4\times4$ Network $\times$ Medium DAG (10 tasks)}

\begin{figure}[h]
\centering
\includegraphics[width=\textwidth]{figures/largenet_mediumdag.png}
\caption{$4\times4$ $\times$ Medium DAG. Widest path funnels all 10 transfers through the
same 8-hop perimeter corridor; shortest path splits traffic across two independent 4--5 hop paths.}
\label{fig:lm}
\end{figure}

Widest path routes flows through longer perimeter corridors, loading links
with multiple concurrent flows. Shortest path uses shorter, more direct routes
with lower average hop counts. Average hops drop from the widest-path route to
shorter corridors, and the makespan improves from 124.86\,s to 101.62\,s
(23\% faster).

%% --- 4x4 Large ---
\clearpage
\subsection{$4\times4$ Network $\times$ Large DAG (20 tasks)}

\begin{figure}[h]
\centering
\includegraphics[width=\textwidth]{figures/largenet_largedag.png}
\caption{$4\times4$ $\times$ Large DAG. Widest path loads perimeter links with up to 12
bidirectional flows; shortest path distributes across the right side of the grid.}
\label{fig:ll}
\end{figure}

Widest path routes transfers through longer perimeter corridors, concentrating
flows on ``best'' links. The topology panel shows thick links on the perimeter
while the right side of the grid remains underused. Shortest path splits
traffic across shorter, more diverse paths with lower average hop counts.
Makespan: 221.73\,s vs.\ 160.62\,s (38\% improvement for shortest path).

%% ============================================================
\section{Discussion}

\subsection{Why Widest Path Fails Under CSMA}

Widest-path routing optimizes a \emph{single-flow} metric: it selects the path
with maximum bottleneck bandwidth. This is optimal when a single transfer uses
the network in isolation. However, DAG execution produces \emph{multiple
concurrent transfers}, and widest-path routing directs them all toward the
same high-bandwidth links. Under CSMA/Bianchi:

\begin{enumerate}
    \item \textbf{Flow concentration}: Multiple transfers share the same
    ``best'' links, dividing effective bandwidth by the number of concurrent
    flows via Bianchi MAC contention ($\eta(n)/n$ scales sublinearly).
    \item \textbf{Hop count inflation}: To reach distant high-bandwidth links,
    widest path takes longer routes. Each additional hop means one more link
    occupied simultaneously, increasing the total contention footprint.
    \item \textbf{Cascading delays}: Longer transfers block downstream tasks,
    pushing the critical path further out and amplifying the makespan.
\end{enumerate}

\subsection{When Widest Path Wins on Grids}

In the grid topology experiments, widest path never wins. The two ties occur
with the small DAG on small networks ($2\times2$ and $3\times3$), where both
strategies use the same single-hop path and the routing decision is irrelevant.
In all non-trivial cases, shortest-path's lower hop count reduces the CSMA
contention footprint enough to outweigh any bandwidth advantage from widest
path.

\subsection{Scaling Behavior}

The widest-path penalty grows with both network size and DAG complexity:

\begin{itemize}
    \item \textbf{Network size}: Larger grids have longer ``widest'' routes
    (up to 8 hops on $4\times4$), amplifying the hop-count penalty.
    \item \textbf{DAG complexity}: More tasks produce more concurrent transfers,
    all competing for the same concentrated corridor.
\end{itemize}

On the $4\times4$ grid, widest path produces makespans 23--338\% longer than
shortest path across all three DAG sizes.

%% ============================================================
\clearpage
\section{Bottleneck Network: When Widest Path Wins}

The grid results suggest shortest-path routing is always superior. However,
this conclusion is topology-dependent. We design a \textbf{bottleneck relay
network} where the shortest path traverses a slow direct link while a
higher-bandwidth relay chain is available.

\subsection{Network Topology}

Five nodes arranged in a line with 30\,m spacing:

\begin{verbatim}
n0(0,0)---n1(30,0)---n2(60,0)---n3(90,0)---n4(120,0)
  \                                          /
   ------------- 120m direct ---------------
\end{verbatim}

All node pairs within 130\,m communication range are connected with
bidirectional links (20 links total). The 802.11ax PHY rates by distance give:

\begin{itemize}
    \item Adjacent links (30\,m): \textbf{8.6\,MB/s} (MCS~5)
    \item 2-hop links (60\,m): 4.3\,MB/s (MCS~3)
    \item 3-hop links (90\,m): 2.15\,MB/s (MCS~1)
    \item Direct n0--n4 link (120\,m): \textbf{1.075\,MB/s} (MCS~0)
\end{itemize}

The key asymmetry: \textbf{shortest path} from n0 to n4 uses the direct
1-hop link at 1.075\,MB/s. \textbf{Widest path} routes through the relay
chain n0$\to$n1$\to$n2$\to$n3$\to$n4 with bottleneck bandwidth 8.6\,MB/s
(an $8\times$ bandwidth advantage), at the cost of 4 hops.

All nodes have identical compute capacity (200 compute\_units/s) so HEFT
cannot avoid cross-network transfers---task placement is forced via
\texttt{pinned\_to} constraints.

\subsection{DAG Structures}

Three DAGs with tasks pinned to force n0$\leftrightarrow$n4 transfers:

\begin{table}[h]
\centering
\caption{Bottleneck DAG configurations. All tasks: compute\_cost=500, all edges: data\_size=10\,MB.}
\label{tab:bn_dags}
\begin{tabular}{lcp{7.5cm}}
\toprule
Label & Tasks & Structure \\
\midrule
Chain       & 5  & T0(n0)$\to$T1(n4)$\to$T2(n0)$\to$T3(n4)$\to$T4(n0). Sequential transfers, no concurrency. \\
Fork-join~3 & 5  & T0(n0)$\to$\{T1(n2),T2(n3),T3(n4)\}$\to$T4(n0). 3 concurrent outbound + 3 concurrent return transfers. \\
Fork-join~8 & 10 & T0(n0)$\to$\{T1--T8 on n1--n4\}$\to$T9(n0). 8 concurrent outbound + 8 concurrent return transfers. \\
\bottomrule
\end{tabular}
\end{table}

\subsection{Results}

\begin{table}[h]
\centering
\caption{Bottleneck network results: widest-path wins all 3 cases.}
\label{tab:bn_results}
\begin{tabular}{l rr rl}
\toprule
DAG & Widest (s) & Shortest (s) & Speedup & Winner \\
\midrule
Chain (5 tasks)       & \textbf{34.74} & 49.71  & $1.4\times$ & Widest \\
Fork-join~3 (5 tasks) & \textbf{40.85} & 101.86 & $2.5\times$ & Widest \\
Fork-join~8 (10 tasks) & \textbf{96.44} & 100.57 & $1.04\times$ & Widest \\
\bottomrule
\end{tabular}
\end{table}

\subsubsection{Chain DAG (sequential transfers)}

With 4 sequential n0$\leftrightarrow$n4 transfers and no concurrency, this is
the cleanest test of raw path bandwidth. Widest path routes each transfer
through the 4-hop relay chain (bottleneck 8.6\,MB/s): each 10\,MB transfer
takes $\sim$1.16\,s plus latency. Shortest path uses the direct 1-hop link
(1.075\,MB/s): each transfer takes $\sim$9.3\,s. The $8\times$ bandwidth
advantage directly translates to a $1.4\times$ makespan improvement (the gap
is less than $8\times$ because compute time dominates).

\subsubsection{Fork-join 3 (moderate concurrency)}

Three concurrent outbound transfers fan out to n2, n3, n4 with varying path
lengths. Widest path routes all three through the relay chain, achieving high
throughput despite CSMA contention among the 4-hop paths. Shortest path sends
all three through their respective direct links (60\,m, 90\,m, 120\,m), but
the n0$\to$n4 transfer at 1.075\,MB/s becomes the bottleneck that determines
when T3 can start. The $2.5\times$ speedup is the largest in the bottleneck
set, demonstrating that the bandwidth disparity overwhelms the contention cost
at moderate concurrency.

\subsubsection{Fork-join 8 (high concurrency)}

With 8 concurrent transfers, the contention cost of 4-hop relay routing
begins to erode the bandwidth advantage. All 8 transfers compete for airtime
on the relay chain's shared links, reducing the per-flow effective bandwidth.
Shortest path also suffers (8 flows share fewer links), but with shorter
paths. The widest-path advantage shrinks to just 4\%, suggesting a crossover
point where high concurrency would make shortest path preferable even on
bottleneck topologies.

\subsection{When Does Widest Path Win?}

The bottleneck experiments identify three conditions for widest-path advantage:

\begin{enumerate}
    \item \textbf{Large bandwidth disparity}: The relay chain must offer
    significantly higher bandwidth than the direct path ($8\times$ here).
    \item \textbf{Low to moderate concurrency}: With sequential or few
    concurrent transfers, the bandwidth gain exceeds the CSMA contention
    penalty. At high concurrency (fork-join~8), the advantage nearly vanishes.
    \item \textbf{Sparse topology}: The network must have a clear bottleneck
    on the shortest path, with no alternative short paths of comparable
    bandwidth.
\end{enumerate}

%% ============================================================
\section{Conclusion}

Under CSMA/Bianchi interference on wireless mesh networks, the optimal routing
strategy depends on network topology:

\begin{itemize}
    \item \textbf{Dense mesh grids}: Shortest-path routing wins 7 of 9
    cases (2 ties). Widest-path's flow concentration and longer routes create
    excessive CSMA contention, with overheads of 23--338\%.
    \item \textbf{Bottleneck relay networks}: Widest-path routing wins all 3
    cases with speedups up to $2.5\times$, leveraging the $8\times$ bandwidth
    advantage of the relay chain over the slow direct link.
\end{itemize}

The key insight is that routing optimality depends on the
\emph{bandwidth-contention tradeoff}: widest-path is beneficial when the
bandwidth gain from relay paths substantially exceeds the CSMA contention cost
of occupying more links, and harmful when dense topologies offer no such
bandwidth advantage but route diversity matters. The fork-join~8 result
(4\% advantage) hints at a crossover point---at sufficiently high
concurrency, even large bandwidth disparity cannot overcome contention costs.
Future work could explore contention-aware routing that jointly optimizes
bandwidth and expected interference load.

\end{document}
